\section{Sharper scalar and planar bounds}\label{sec:geometry}
The coarse algebraic region bound above works in every fixed dimension and
allows atoms. The scalar and planar estimates here give sharper dependence
on the smoothing scale under continuous noise. We then transfer formula
counts to finite grids. Throughout, the coordinate penalties are independent;
branch costs themselves need not be independent.

\subsection{Scalar ties and formula intervals}
\begin{theorem}[Scalar continuous bound]\label{thm:scalar}
Let the deterministic branches $q_z$, $z\in Z\subseteq\{0,1\}^m$, be
$L$-Lipschitz on a compact interval of length $T>0$. If the $\xi_i$ have
densities at most $\phi$, the number $H$ of interior points with at least
two minimizing labels and the number $I$ of connected open intervals of
unique minimizing support satisfy
\[
 \E H\le2m\phi LT,\qquad \E I\le1+2m\phi LT.
\]
These statements include isolated touching ties and do not require polynomial
branches.
\end{theorem}
\begin{proof}
At a fixed parameter $t$, let $\Delta(t)$ be the difference between the two
smallest distinct-label values, counting ties as zero and using infinity for
a one-label family. For each coordinate $i$, condition on all other noise.
The two bit-class minima have the form $a_i$ and $b_i+\xi_i$, with
$a_i,b_i$ fixed. If the two best labels differ on $i$, they attain these
two class minima. The conditioning argument of Beier and V\"ocking's generalized
isolating lemma \cite[conference version, Lemma~5]{BeierVocking} therefore gives
\begin{equation}\label{eq:scalargap}
 \Prob\{\Delta(t)\le d\}\le2m\phi d\qquad(d\ge0).
\end{equation}
Empty bit classes contribute nothing.

Partition the parameter interval into $r$ equal cells of length $\eta=T/r$.
Almost surely, every endpoint and midpoint in these countably many meshes
has a unique winner. If a cell contains a tie at $t$ and $z$ wins at its
midpoint $s$, at least one tied label $y$ differs from $z$. Then
\[
 0\le F_y(s)-F_z(s)\le2L|s-t|\le L\eta.
\]
Thus the number $C_r$ of cells containing a tie has
$\E C_r\le2m\phi LT$ by \eqref{eq:scalargap}. Every finite subset of
the tie set occupies distinct cells for all sufficiently fine meshes.
It follows that $H\le\liminf_r C_r$, also when $H$ is infinite. Fatou's
lemma proves the first bound, and hence almost-sure finiteness. Between tie
points the unique winning label is locally constant, and therefore constant
on each interval. Endpoint ties have probability zero, giving $I\le H+1$.
The events are measurable: the gap is continuous in the finite branch-value
vector and its minimum over a compact cell is continuous in the noise.
\end{proof}

For quadratic branches define a \emph{formula interval} as a maximal
positive-length interval on which the envelope has one fixed polynomial
formula. Identify identical polynomials, and do not split an interval at an
isolated touching tie when the envelope formula stays the same. Separated
appearances of one formula count separately.

\begin{lemma}[Full univariate quadratic envelopes]\label{lem:DS}
The lower envelope of $r\ge1$ distinct full univariate quadratics has at most
$2r-1$ formula intervals. It can be constructed in $O(r\log(r+1))$
arithmetic, comparison, and quadratic-root operations.
\end{lemma}
\begin{proof}
The sequence of interval labels has no adjacent repetition and no subsequence
$a,b,a,b$: four strictly winning witness points would force at least three
roots of the nonzero quadratic difference. A sequence on $r$ symbols with
these two properties has length at most $2r-1$. Indeed let its first symbol
$a$ occur $s$ times. The blocks after successive occurrences of $a$ have
disjoint alphabets, since a shared symbol would give $a,b,a,b$. At least
$s-1$ blocks are nonempty. Applying induction within each nonempty block
gives total length at most $s+2(r-1)-(s-1)=2r-1$.
This is the elementary order-two Davenport--Schinzel bound.
To merge two envelopes, traverse their overlapping intervals. On each overlap
the two current quadratics have at most two crossings, so merging is linear
in the two input lengths. Divide and conquer, using the proved linear
length bound at every level, gives the operation count. These are classical
envelope operations, not new algorithmic primitives.
\end{proof}

\begin{theorem}[Scalar atomic bound]\label{thm:scalargrid}
For quadratic $L$-Lipschitz branches on an interval of length $T$, and the
grid noise \eqref{eq:noisegrid}, the formula-interval count $I_{\rm f}$ obeys
\begin{equation}\label{eq:scalargrid}
 \E I_{\rm f}\le1+2m\phi LT+\frac{4m2^m}{N},
 \qquad \phi=\frac1{2\sigma}.
\end{equation}
\end{theorem}
\begin{proof}
Add independent continuous jitter: $\eta_i=\xi_i+\delta U_i$, with
$U_i$ uniform on $[-1,1]$ and $\delta>0$. Conditioning on $U_i$ and
translating an interval in \eqref{eq:gridmass} shows that
\[
 \Prob\{\eta_i\in J\}\le\phi|J|+1/N.
\]
The variables $\eta_i$ are independent and have densities; the displayed
bound is uniform in $\delta$.

Condition on all but $\eta_i$ and define
\[
 h_i(t)=\min_{z_i=0}\left(q_z(t)+\sum_{j\ne i}\eta_jz_j\right)
       -\min_{z_i=1}\left(q_z(t)+\sum_{j\ne i}\eta_jz_j\right),
\]
ignoring an empty class. Each class has at most $2^{m-1}$ branches.
Lemma~\ref{lem:DS} bounds the common refinement of the two envelopes by
fewer than $2^{m+1}$ pieces. Splitting each quadratic piece at its stationary
point gives at most $2^{m+2}$ monotone arcs. The function $h_i$ is
$2L$-Lipschitz, so its total variation is at most $2LT$. Constant-arc
values and arc endpoint values have conditional probability zero for
$\eta_i$. Every remaining arc contributes at most one level hit and its
image has probability at most its length times $\phi$, plus $1/N$.
Therefore
\[
 \E\#\{t:h_i(t)=\eta_i\}\le2\phi LT+2^{m+2}/N.
\]
Any optimal tie differs on some bit, whose class minima coincide. Summing
over bits bounds the jittered tie count, and hence its unique-support
interval count, by the right side of \eqref{eq:scalargrid}.

For a fixed unjittered outcome choose ordered interior witness points, one
per formula interval, avoiding the finitely many zeros of nonidentical
polynomial differences. At each point the winning polynomial has a positive
gap from every different polynomial. For all sufficiently small $\delta$,
every jittered winner there belongs to this original polynomial class.
Consecutive witnesses belong to different classes and require distinct
jittered intervals. Along a fixed sequence $\delta\downarrow0$, ties at
these points have probability zero simultaneously. Thus the unjittered
formula-interval count is at most the liminf of the jittered counts.
Fatou's lemma proves \eqref{eq:scalargrid}. The one-label case $m=0$
needs no jitter.
\end{proof}

\subsection{Compact planar levels and their curvature}
The geometric mechanism below is classical. Bourgain, Korobkov, and
Kristensen \cite[Theorem~6.1, Corollary~6.2, and their proof]{BourgainLevels}
use Hessian variation, coarea, and total curvature to control planar
bounded-Hessian levels; Ambrosio and Bertrand \cite[Section~5.1]{AmbrosioDC}
develop measure-valued Hessians in a broader difference-of-convex calculus.
We give the finite piecewise-quadratic argument, including its interface
terms, to specify the constant needed here.

\begin{lemma}[Compact-level bound]\label{lem:loops}
Let $D=[-M,M]^2$, $M>0$, with perimeter $P=8M$. Suppose $g_0,g_1$
are continuous concave $L$-Lipschitz functions, each quadratic on a finite
semialgebraic partition of $D$. For almost every $y$, the level set of
$h=g_0-g_1$ in $\operatorname{int}D$ is a finite union of regular arcs
with finitely many corners. If $C(y)$ counts its compact connected
components, then
\begin{equation}\label{eq:loopintegral}
 \int_\R C(y)\,dy\le LP/\pi.
\end{equation}
\end{lemma}
\begin{proof}
Refine the partitions to a finite smooth semialgebraic stratification of
cells, interfaces, and vertices, including the boundary. Exclude the critical
values of the restriction of $h$ to every stratum. There are only finitely
many such values: each critical locus is semialgebraic with finitely many
connected components, and the function is constant on each component.
Constant restrictions contribute their constant value. Every remaining
level is regular within cells and transverse to interfaces, avoids vertices,
and consists of finitely many piecewise smooth curves.

Write $|D^2h|_*$ for the total variation of its matrix-valued distributional
Hessian using the nuclear norm (the sum of singular values). On a smooth
cell, if $v$ is the unit tangent of a regular level curve, its absolute
curvature is $|v^T(D^2h)v|/\|\nabla h\|$. The coarea formula therefore
bounds the integral over levels of smooth-arc curvature by
$\int\|D^2h\|_*$. To include the corners, choose tangent and normal
coordinates at an interface. Continuity gives one-sided gradients
$(a,b_-)$ and $(a,b_+)$ with the same tangential derivative. Transversality
means $a\ne0$. The angle $\alpha$ between the two level tangents satisfies
\begin{equation}\label{eq:cornercharge}
 |a|\alpha\le|b_+-b_-|.
\end{equation}
Indeed the gradients lie in the same half-plane with fixed nonzero tangent
coordinate $a$; the derivative of their angle as a function of $b$ has
absolute value at most $1/|a|$. Their tangent lines turn through this same
angle. One-dimensional coarea along the interface multiplies its corner
angle by $|a|$, and \eqref{eq:cornercharge} bounds the resulting integral
by $\int|b_+-b_-|\,ds$. This is exactly the nuclear-norm mass of the
singular Hessian there: the gradient jump is normal, so its Hessian density
is a rank-one matrix with that magnitude. Vertices have been excluded by
their values. Thus, with total curvature including corners,
\begin{equation}\label{eq:curvaturecoarea}
 \int_\R\operatorname{TC}(\{h=y\}\cap\operatorname{int}D)\,dy
 \le |D^2h|_*(\operatorname{int}D).
\end{equation}

For concave $g$, the measure $-D^2g$ is positive semidefinite. Its nuclear
mass equals its negative trace, $-\Delta g$, and is at most $LP$.
To see the last inequality explicitly, use inner parallel squares with
generic boundaries. The piecewise smooth divergence theorem, including
the gradient jumps, identifies the trace mass with minus the outward
normal gradient flux. Its absolute value is at most $L$ times that square's
perimeter. Increasing the squares to the interior of $D$ and taking the
monotone limit proves the bound. Consequently
$|D^2h|_*\le |D^2g_0|_*+|D^2g_1|_*\le2LP$.
Every compact regular component is a simple closed curve: local regularity
and transverse interface crossings give a one-dimensional manifold without
boundary. Its total absolute curvature is at least $2\pi$, including
corners, by Fenchel's theorem \cite[Theorem~2.4]{SullivanCurvature}.
Divide \eqref{eq:curvaturecoarea} by $2\pi$ to obtain
\eqref{eq:loopintegral}.
\end{proof}

\subsection{A planar continuous region theorem}
\begin{theorem}\label{thm:planar}
Let $q_z$, $z\in Z\subseteq\{0,1\}^m$, be quadratics on
$D=[-M,M]^2$, $M>0$. Suppose one common quadratic $p$ makes every
$r_z=q_z-p$ concave with $\|\nabla r_z\|\le L$ on $D$.
For independent noises with densities at most $\phi$, let $\mathcal R$
count connected open regions in $\operatorname{int}D$ with a unique
minimizing label. With $A=4M^2$ and $P=8M$,
\begin{equation}\label{eq:planar}
 \E\mathcal R\le1+(1+1/\pi)m\phi LP
                  +96\binom{m}{2}\phi^2L^2A.
\end{equation}
\end{theorem}
\begin{proof}
Subtract $p$. For a coordinate $i$ with two nonempty bit classes, condition
on all other noise and put
\[
 g_a(t)=\min_{z_i=a}\left(r_z(t)+\sum_{j\ne i}\xi_jz_j\right),
 \qquad h_i=g_0-g_1.
\]
Each $g_a$ is concave and $L$-Lipschitz; hence $h_i$ is $2L$-Lipschitz.
The equality $h_i(t)=\xi_i$ means exactly that both classes attain the
global minimum. Along each boundary edge, one-dimensional coarea and the
conditional density bound give at most $2\phi L$ times the edge length
expected hits. Every optimal tie is witnessed by a differing bit, so the
number $B$ of boundary tie points satisfies
\begin{equation}\label{eq:planarB}
 \E B\le2m\phi LP.
\end{equation}
Corners are almost surely untied. Quadratic differences restricted to edges
have finitely many critical values (or are constant), so almost surely
every boundary pair equality is transverse.

Let $J$ count interior points with at least three minimizing labels.
Three distinct binary vectors have affine rank two, and some coordinate
pair $i,j$ projects them to three corners of $\{0,1\}^2$. Condition
outside that pair and define its four class minima $C_{ab}(t)$, omitting
empty classes. Each is $L$-Lipschitz. A corner triple has an elbow and two
neighbors; equality of their three class values gives
$(\xi_i,\xi_j)=H(t)$, where each component of $H$ is a signed difference
of two class minima. Its derivative rows have norms at most $2L$ almost
everywhere, so $|\det DH|\le4L^2$. The area formula and conditional density
at most $\phi^2$ bound the expected number of its interior solutions by
$4\phi^2L^2A$. There are four corner triples per coordinate pair. Therefore
\begin{equation}\label{eq:planarJ}
 \E J\le16\binom{m}{2}\phi^2L^2A.
\end{equation}
Class-equality solutions not globally optimal merely overcount. The image
under $H$ of a boundary edge has planar measure zero, so there are almost
surely no boundary triple ties.

Let $W$ count compact components of the global tie set in
$\operatorname{int}D$ with no triple junction; boundary-to-boundary arcs
are excluded. On each such loop the same two labels minimize throughout;
a change would create a triple tie. Choose a bit on which those labels
differ. The loop is an entire compact component of $\{h_i=\xi_i\}$:
a further branch of that level would also be a global tie and would require
a triple junction. Lemma~\ref{lem:loops}, applied conditionally, gives
\begin{equation}\label{eq:planarW}
 \E W\le m\phi LP/\pi.
\end{equation}

It remains to bound the graph formed by the ties. Almost surely every
pairwise quadratic zero curve is regular or empty. Indeed its nonconstant
part is fixed, whereas its constant shift has a density; a singular level
of a quadratic must be its unique critical value, if it has one. A constant
difference is almost surely nonzero. Also, almost surely no five labels
co-minimize anywhere. Five cube points have affine rank at least three,
as in Proposition~\ref{prop:rankgap}. Choose four with independent
differences and three noise coordinates giving an invertible difference
matrix. Conditional on the other noise, their equalities place those three
coordinates in the image of a polynomial map from $D$ to $\R^3$. That
image has three-dimensional measure zero. There are only finitely many
choices, proving the assertion. Four-way ties are not excluded.

Away from $J$ and $B$, a tie is locally a smooth arc between two labels;
there is no isolated two-label tie because its equality curve is regular
and all other labels remain strictly higher locally. At a junction at most
four labels meet, giving at most six pairwise curves and hence degree at
most twelve. Semialgebraicity and \eqref{eq:planarJ} give a finite graph
almost surely. After separating its junction-free loops, the remaining
arcs have at most $12J+B$ ends, counted with multiplicity, and hence number
at most $6J+B/2$. Isolated junctions do not separate the square. Adding an
arc or closed loop increases the number of complementary regions by at most
one; equivalently, Euler's formula gives
\[
 \mathcal R\le1+6J+B/2+W.
\]
Every complementary region has a fixed unique winner. Substitution of
\eqref{eq:planarB}--\eqref{eq:planarW} proves \eqref{eq:planar}.
\end{proof}

Both qualifications in the graph argument matter. The affine binary family
$F_z(t)=z_1(t_1+\xi_1)+z_2(t_2+\xi_2)$ has a four-way tie when
$(-\xi_1,-\xi_2)$ lies inside the square. The pair $0$ and
$\|t\|^2-r^2$, $0<r<M$, has a compact tie loop with no junction;
small constant perturbations preserve such a loop. The term $W$ cannot
be omitted.

The normalization is available for every family of quadratics whose gradients
have norm at most $K_0$ on $D$. For a unit vector $u$, evaluate the affine
gradient at $Mu$ and $-Mu$ to obtain
$\|\nabla^2q_z\|_{\rm op}\le K_0/M$. Thus
$p(t)=K_0\|t\|^2/(2M)$ makes $q_z-p$ concave, with gradient bound
$L=(1+\sqrt2)K_0$. More directly, the boundary-excluded message branches
\eqref{eq:branch} are already concave. Theorem~\ref{thm:planar} applies
to a two-dimensional spectral message on $[-R,R]^2$, $R>0$, by taking
the planar half-width $M=R$ and $L=L_2$ from \eqref{eq:boundL}, without
diagonal dominance.

\subsection{Distribution discrepancy and necessary formulas}
For polynomial branches on a nonempty open domain $T\subseteq\R^k$,
let $U(\eta)$ count labels uniquely minimizing somewhere on $T$.
Let $K_{\rm f}(\eta)$ count distinct polynomial formulas equal to the
envelope on a nonempty open subset, identifying identical polynomials.
These are the necessary formulas. They still represent the envelope everywhere:
outside the finite union of zero sets of nonidentical differences, a winning
formula is unique; that set is dense, and continuity supplies the remaining
points. On a box, the same reasoning extends coverage to its closed boundary.

\begin{theorem}[Monotone discrepancy transfer]\label{thm:transfer}
Let $\mu=\prod_i\mu_i$ and $\nu=\prod_i\nu_i$. Suppose each pair of
marginals differs by at most $\Delta_i$ on every half-line, with either
open or closed endpoint convention. Then
\begin{equation}\label{eq:Utransfer}
 |\E_\nu U-\E_\mu U|\le |Z|\sum_i\Delta_i.
\end{equation}
If $\E_\mu U_{q+c}\le B$ uniformly over deterministic constant branch
shifts $c$, then
\begin{equation}\label{eq:Ktransfer}
 \E_\nu K_{\rm f}\le B+|Z|\sum_i\Delta_i.
\end{equation}
If instead, for a positive integer $p$, the uniform bound is
$\E_\mu U_{q+c}^p\le B_p$, then
\begin{equation}\label{eq:transfermoments}
 \E_\nu K_{\rm f}^p\le B_p+2|Z|^p\sum_i\Delta_i.
\end{equation}
\end{theorem}
\begin{proof}
Fix $z$ and all noise coordinates except $i$. If $z_i=0$, its unique
activity persists as $\eta_i$ increases: same-bit comparisons stay fixed
and opposite-bit comparisons improve. Thus the section of its activity
event is an upward half-line, possibly empty or the whole line. For
$z_i=1$ it is a downward half-line. Activity is measurable, since a strict
witness can be chosen from a countable dense subset of $T$. Replace the
marginals one at a time, condition on the others, and telescope. The change
in the probability of this label's activity is at most $\sum_i\Delta_i$.
Summing over labels proves \eqref{eq:Utransfer}. This argument uses only
continuous branches.

To handle persistent duplicate polynomials, choose distinct fixed numbers
$a_z$ and consider $q_z+\delta a_z$ as $\delta\downarrow0$. For each
necessary original formula, choose a witness in its open winning set outside
all zero sets of nonidentical differences. It has a positive gap from every
different polynomial. In its identical-formula class, the label with smallest
$a_z$ becomes the unique winner at that witness for all sufficiently small
$\delta>0$. There are finitely many classes, so, pointwise in $\eta$,
\[
 K_{\rm f}(\eta)\le\liminf_{\delta\downarrow0}U_{q+\delta a}(\eta).
\]
Fatou's lemma along $\delta=1/j$ and \eqref{eq:Utransfer} prove
\eqref{eq:Ktransfer}. The shifts are a proof device only.

For the last assertion expand $U^p$ over ordered $p$-tuples of labels,
allowing repetitions. A coordinate section of their simultaneous activity
is an intersection of upward and downward half-lines, hence an interval.
Its discrepancy is at most $2\Delta_i$. The same marginal replacement
bounds the difference of the $p$th moments by
$2|Z|^p\sum_i\Delta_i$. Apply the preceding pointwise limit, raised to
$p$, and Fatou's lemma to obtain \eqref{eq:transfermoments}.
\end{proof}

For uniform continuous noise on $[-\sigma,\sigma]$ and its
endpoint-inclusive $N$-point grid, $\Delta_i\le1/N$. After rescaling to
$[0,1]$, the grid distribution function is constant between points
$j/(N-1)$; comparison with the identity at both ends of every step proves
the bound, including open endpoints. Consequently
\[
 \E_{\rm grid}K_{\rm f}\le B+m2^m/N.
\]
Constant shifts preserve the assumptions of Theorems~\ref{thm:scalar}
and \ref{thm:planar}, and $U$ is bounded by their region counts. Thus one
may use either
\[
 B=1+2m\phi LT
 \quad\text{or}\quad
 B=1+(1+1/\pi)m\phi LP+96\binom{m}{2}\phi^2L^2A.
\]
For transfer error at most $\epsilon>0$, a power-of-two grid with
$N\ge\max\{2,m2^m/\epsilon\}$ uses
$O(m+\log(m+1)+\log_+(1/\epsilon))$ bits per coordinate. Its cardinality
is exponential, although its samples have polynomial bit length.
These sharper constants concern distinct necessary formulas, not all active
labels or connected regions. In particular the scalar correction
$m2^m/N$ does not replace $4m2^m/N$ in the formula-interval theorem.
The integer-moment transfer has error at most $2m2^{mp}/N$; it requires
a continuous $p$th-moment bound, which the planar first-moment theorem
alone does not supply. The direct moment bounds in
Appendix~\ref{sec:moments} use much smaller grids when sharp first-moment
constants are unnecessary.
